\documentclass[11pt]{amsart}
\usepackage[margin=1.2in]{geometry}
\usepackage{amsmath,amssymb}
\usepackage{booktabs}
\usepackage[hidelinks]{hyperref}
\usepackage{url}

\newtheorem{theorem}{Theorem}
\newtheorem{lemma}[theorem]{Lemma}

\newcommand{\tdiv}{\tau}
\newcommand{\Nat}{\mathbb{N}}

\title{A kernel-checked exclusion certificate for Erd\H{o}s problem 647}

\author{Ibby Mian}
\author{Shayaan Siddique}
\address{Millennium Research}
\email{ibrahimnmian@gmail.com}

\date{August 2026}

\begin{document}

\begin{abstract}
Erd\H{o}s problem 647 asks whether any $n > 24$ satisfies
$\max_{m<n}(m + \tdiv(m)) \le n + 2$, where $\tdiv$ is the divisor-count
function. Computational searches have excluded solutions up to $10^{12}$
by direct sieve (Id\'en) and up to roughly $9.17 \times 10^{18}$ within a
modular reduction whose Lean component relies on \texttt{native\_decide}
(Hughes, extended by bentrd). Those computations sit outside any proof
kernel. We give the first exclusion checked end to end by one: no solution
exists with $24 < n \le 10^8$, proved in Lean~4 with axiom closure exactly
$\{\mathtt{propext}, \mathtt{Classical.choice}, \mathtt{Quot.sound}\}$.
The proof replays a chain of $891{,}554$ factorization witnesses whose
excluded intervals concatenate across $(24, 10^8]$; it needs no primality
facts beyond primes below $1024$, and it is the finite, fully proved form
of a domination-interval argument whose asymptotic step was the identified
gap in a withdrawn January 2026 claim on this problem. Our range is four
to eleven orders of magnitude below the computational frontiers we cite; the
contribution is the trust base, not the range.
\end{abstract}

\maketitle

\section{The problem and its status}

Let $\tdiv(m)$ denote the number of divisors of $m$. Problem 647 in
Bloom's database of Erd\H{o}s problems \cite{bloom647}, due to Erd\H{o}s
and Selfridge, asks:
\begin{quote}
Is there some $n > 24$ such that
$\max_{m<n}\,(m + \tdiv(m)) \le n + 2$?
\end{quote}
The condition holds at $n = 24$, and the threshold $n+2$ is best possible,
since $\max(\tdiv(n-1)+n-1,\ \tdiv(n-2)+n-2) \ge n+2$ for every $n$.
Erd\H{o}s considered it ``extremely doubtful'' that infinitely many such
$n$ exist and suggested that $\max_{m<n}(\tdiv(m)+m-n) \to \infty$
\cite{erdos79}, while expecting the windowed variant
($\max_{n-k<m<n}$ for fixed $k$) to have infinitely many solutions for
every $k$; in 1992 he offered \pounds 25 for a single example beyond
$24$. The problem page records the further original references and
Tao's observation that the problem is comparable to the first part of
problem 679 and much stronger than problems 413 and 248
\cite{bloom647}. Guy's account of the neighboring questions is
\cite[B8]{guy}. The known solutions of the associated equality,
OEIS A087280, are $5, 8, 10, 12, 24$ \cite{oeisA087280}.

The interval $(24, X]$ contains no solution for increasingly large $X$
by computation. McCranie's bound $X = 10^{10}$ stood in OEIS A087280
until June 2026, when Id\'en ran a segmented multiplicative sieve over
all $n \le 10^{12}$, tracking the running maximum
$R(n) = \max_{m<n}(m + \tdiv(m))$, and found no solution in
$(24, 10^{12}]$; the minimum observed gap $R(n) - n$ was $224$, near
$n = 10^{11}$ \cite{iden}. In parallel, a sequence of structural results
accumulated in the problem's public discussion thread. Dutta showed that
any solution $n > 84$ satisfies $2520 \mid n$ and, writing $n = 2520N$,
extracted primality constraints from the divisor bound at small shifts;
Alexeev refined the $2$-adic bookkeeping, upgraded $2520N - 1$ to a
prime, and exhibited the near-solution $n = 604{,}517{,}614{,}941{,}240$,
which satisfies the divisor bound at every shift $k \le 13$ and at
$k = 24$. Kitamura added that $(n-3)/3$ must be prime. Hughes organized
these into a sieve modulo $11 \cdot 13 \cdot 17 \cdot 19 = 46189$ leaving
$41$ open residue classes, showed every solution beyond $24$ lies in one
of two admissible prime $7$-tuple families, and derived the unconditional
bound $|\mathcal{C}(x)| \ll x/(\log x)^7$ for the count of candidates up
to $x$ by Brun's sieve \cite{hughesrepo}. Hughes also published a
\emph{frontier certificate}: no solution with
$24 < n \le 2520 \cdot 46189 \cdot 529 \cdot 10^7 \approx 6.16 \times
10^{17}$, obtained by closing $6{,}549$ refined sub-progressions with
fixed congruences and searching the remaining $15{,}140$ to depth
$u = 10^7$; bentrd later extended the search leg to
$u = 1.49 \times 10^8$, i.e.\ to $n \le 9.17 \times 10^{18}$
\cite{bentrd}. The Lean component of Hughes's development is explicit
that its finite computations are discharged by \texttt{native\_decide}
and that one problem-specific axiom carries the open residue classes
\cite{hughesrepo}.

One further episode frames our contribution. In January 2026 a claimed
full solution appeared \cite{agbanwa}: a \emph{domination} argument in
which each record-holder of $m \mapsto m + \tdiv(m)$ excludes an explicit
interval of $n$, together with the assertion that these intervals cover
all $n > 24$. Tao's public assessment identified the gap precisely: the
claim that consecutive domination intervals overlap forever is unproven
and likely false, and the accompanying Lean formalization treats the
asymptotic case as an axiom, so it certifies nothing \cite{bloom647}.
The finite form of the same idea, however, is exactly what a proof
kernel can settle: over a bounded range, whether the intervals cover is
a checkable fact.

This note reports that check, carried out at scale under a strict trust
discipline. Our result is:

\begin{theorem}\label{thm:main}
No $n$ with $24 < n \le 10^8$ satisfies
$\max_{m<n}(m + \tdiv(m)) \le n + 2$.
\end{theorem}

Theorem~\ref{thm:main} is proved in Lean~4 \cite{lean4} against
mathlib \cite{mathlib}, in the repository \texttt{decanus}
\cite{decanus}. The axiom closure of the final theorem, reported by
\verb|#print axioms| and enforced in continuous integration by a
two-layer gate, is exactly the three standard axioms
$\{\mathtt{propext}, \mathtt{Classical.choice}, \mathtt{Quot.sound}\}$:
no \texttt{sorry}, no \texttt{native\_decide}, no problem-specific
axiom, no external solver or compiled evaluator in the trusted base.
The statement is formulated against the inner predicate of the
\texttt{formal-conjectures} formalization of the problem
\cite{formalconjectures}, pinned at a fixed commit, so the theorem
speaks the same vocabulary as the public formal statement. A companion
theorem certifies the subrange $(24, 10^7]$ separately.

We state the position of this result plainly, in the same spirit as our
earlier certifications of Erd\H{o}s problems 7, 364, and 486. It is not
a computational record: Id\'en's $10^{12}$ and the
Hughes--bentrd frontier at $9.17 \times 10^{18}$ are far larger, and we
cite them as the state of the search. What the present work adds is a
range on which nothing needs to be trusted beyond the Lean kernel and
its three standard axioms --- not contributor code, not a compiler, not
a cache, not a run of a sieve on one machine. Section~\ref{sec:cert}
describes the certificate, Section~\ref{sec:lean} the Lean development,
Section~\ref{sec:validation} the generation pipeline and its
cross-validation, and Section~\ref{sec:limits} the method's ceiling and
its relation to the structural program.

\section{The certificate}\label{sec:cert}

Call $m$ a \emph{kill witness} for $n$ if $m < n$ and
$m + \tdiv(m) \ge n + 3$; one such witness refutes the defining
inequality at $n$. A single $m$ serves every
$n \in [m+1,\, m + \tdiv(m) - 3]$, so a finite set of witnesses whose
intervals concatenate without gaps excludes an entire range. The
certificate for Theorem~\ref{thm:main} is such a chain: witnesses
$m_1, m_2, \ldots$ with associated certified bounds $t_i \le \tdiv(m_i)$,
each interval $[m_i + 1,\, m_i + t_i - 3]$ beginning no later than the
previous one ends, jointly covering $(24, 10^8]$.

Two design choices make the chain cheap for a kernel to check.

\emph{Lower bounds only, from maximal smooth parts.} A witness never
needs its full factorization certified; it needs a lower bound on
$\tdiv(m)$. Each witness stores the maximal power of every prime
$p < 1024$ dividing $m$. The checker re-multiplies the stored powers,
verifies that their product divides $m$ and that each listed power is
maximal (i.e.\ $p^{e+1} \nmid m$), and reads off
$\prod_i (e_i + 1) \le \tdiv(m)$, since the stored prime powers are
pairwise coprime. Maximality buys one further factor: whatever cofactor
of $m$ remains above the smooth part is then coprime to it, so if the
cofactor exceeds $1$ it contributes at least the two divisors $1$ and
itself, and the certified bound doubles. Primality of a listed $p < 1024$
is settled by eleven trial divisions (by the primes up to $31$, since
$32^2 = 1024$), so no primality certificate for a large prime appears
anywhere in the development.

\emph{Greedy chaining.} Given the certified bound function, the shortest
chain is produced greedily: from cover position $C$ (initially $24$),
take a witness $m \le C$ maximizing $m + t(m)$ and advance to
$m + t(m) - 3$. Greedy choice is optimal for interval covering, and the
resulting chain is the iteration of the running-maximum function
familiar from the sieve computations. Measured chain lengths are
$495$ to cover $(24, 10^4]$, $2{,}829$ to $10^5$, $17{,}941$ to $10^6$,
$123{,}323$ to $10^7$, and $891{,}554$ to $10^8$ --- growth by a factor
of roughly $6.5$--$7$ per decade, reflecting the slow growth of typical
divisor sums. With the cofactor doubling in place the certified chain
for $(24, 10^7]$ has exactly the length of the chain computed from the
true divisor function, an equality we verified directly: the greedy
steps are dominated by record-setters of the form
(smooth part) $\times$ (single large prime), for which the doubled
smooth bound equals $\tdiv$ exactly.

The doubling rule was forced by data, and the episode is worth
recording. A first design certified only the smooth part's
$\prod (e_i + 1)$. The generator, which is required to fail loudly if
the chain ever stalls, stopped at cover position $69{,}030{,}671$: the
best available witness was
$m = 69{,}030{,}650 = 2 \cdot 5^2 \cdot 13 \cdot 61 \cdot 1741$, whose
true divisor count $48$ is halved to $24$ by the prime $1741 > 1024$,
and at this gap minimum the halved bound advanced the cover by zero.
Requiring maximality of the stored powers, which makes the cofactor
coprime for free, recovered the factor of two and unstuck the chain.
Deep gap minima are precisely where solutions would live, so it is not
an accident that this is where a weakened bound first failed.

\section{The Lean development}\label{sec:lean}

The development separates a small trusted-once soundness layer from
bulk data that the kernel evaluates.

The data layer encodes a witness as an offset from the current cover
position together with its list of prime-power pairs, most significant
prime first. A Boolean checker \texttt{chainOk} folds down the witness
list: for each witness it verifies the offset is admissible, the listed
primes are strictly decreasing and pass the eleven-division primality
test, the smooth product divides $m$, and each stored power is maximal
in $m$; it then advances the cover by the certified bound and finally
demands the target is reached. All recursion is structural, with no
\texttt{Nat.sqrt}, no well-founded recursion, and no \texttt{Finset}
computation in the evaluated path, so kernel reduction is
GMP-arithmetic on literals plus list traversal.

The soundness layer proves, once and abstractly, that a passing check
means what it should: that the primality test below $1024$ is sound
(via the bound $\mathrm{minFac}(p)^2 \le p$ for composite $p$); that
the divisor count of the smooth part is exactly $\prod (e_i+1)$
(multiplicativity over the pairwise-coprime stored powers); that
maximality forces the cofactor coprime, giving the doubling; that
divisor counts are monotone under divisibility; and that a successful
\texttt{chainOk} run yields, for every $n$ in the covered interval, a
witness $m < n$ with $n + 3 \le m + \tdiv(m)$. A bridge lemma converts
that witness into the negation of the formalized inequality
$\bigl(\sup_{m : \mathrm{Fin}\, n}\, m + \sigma_0\, m\bigr) \le n + 2$,
the inner predicate of the \texttt{formal-conjectures} statement of the
problem, which the repository pins at a fixed upstream commit and
carries verbatim up to one coercion the elaborator inserts in both
formulations \cite{formalconjectures}.

The certificate data is generated into $254$ chunk files ($31$ for the
$10^7$ rung, $223$ for the $10^8$ rung, about $42$~MB of Lean source in
total), each containing one witness-list literal and one theorem proved
by a single \texttt{decide}: that \texttt{chainOk}, started at that
chunk's opening cover position, reaches its closing one. Elaborating a
chunk takes roughly $40$ seconds, most of it the kernel's own replay of
the evaluation. Generated driver files compose the chunk theorems into
the rung theorems; the composition is forced by the typechecker, since
each chunk's closing position is the next chunk's literal opening
position, and a gap cannot elaborate.

Trust is enforced mechanically. A curated manifest lists the published
theorems and prints their axiom closures; independently, an audit file
walks every theorem of every module of the development in the compiled
environment --- $579$ of them --- recomputes each axiom closure, and
fails to compile if anything exceeds the three standard axioms, if any
\texttt{sorry} survives, or if any native-code axiom appears.
Continuous integration runs both layers on every push. The two-rung
build compiles in Lean 4.30.0 against mathlib v4.30.0 in about three
CPU-hours.

\section{Generation and cross-validation}\label{sec:validation}

The chain is found by a segmented sieve over certified divisor bounds:
an array pass accumulates $\prod (e_i+1)$ over prime powers below
$1024$ and, in the same pass, the smooth part itself, marking entries
whose smooth part falls short of the number --- exactly the entries
whose certified bound doubles. A greedy walk over the running maximum
emits the witness ledger in an append-only JSONL file whose header and
footer record the exact covered interval. Generation of the $10^8$
chain takes about a minute.

Because the generator is the one component a kernel cannot vouch for,
it is not trusted. A second implementation, sharing no code with the
first --- pure Python, trial division throughout --- replays the ledger
witness by witness: primality of every listed prime, strict ordering,
divisibility, maximality of every power, the doubling condition,
strictly increasing cover, and the final bound. A third implementation
recomputes the divisor function by an independent sieve and confirms
that the solution set below $10^7$ is exactly
$\{2,3,4,5,6,8,10,12,24\}$, in agreement with OEIS A087280 and its
conventions, and that $n + \tdiv(n)$ matches OEIS A062249 on the
recorded values \cite{oeisA062249}. Regeneration is deterministic:
re-running the pipeline reproduces the committed chunk files byte for
byte, a property checked by diff in our verification runs. In the other
direction, any dishonesty in the scripts is caught downstream, since
the kernel replays every witness from scratch; the scripts affect
completeness of the search for a chain, never soundness of the theorem.

Beyond the kernel replay that every build performs, we re-checked the
compiled development with the standalone checker
\texttt{lean4checker} \cite{lean4checker}, which re-typechecks the
compiled environment outside the elaborator.
% TO FILL BEFORE SUBMISSION: lean4checker verdict line
% (replay of all 262 modules of the development on the build machine).
A from-source verification leg on separate x86 hardware --- Lean
toolchain compiled from source, mathlib and the development rebuilt
with no cache, certificates regenerated and diffed against the
committed files, gate re-run, environment re-checked ---
% TO FILL BEFORE SUBMISSION: pod leg results and olean digest comparison.
follows the protocol of our Erd\H{o}s 486 audit.

\section{Limits and prospects}\label{sec:limits}

The method's ceiling is set by certificate size and kernel time, not by
mathematics. Chain length grows by a factor of about $6.5$--$7$ per
decade of range; at $10^9$ the certificate would run to roughly
$6$ million witnesses and $260$~MB of source, past what a repository
should carry, and at $10^{12}$ --- the sieve frontier --- to around
$2 \times 10^9$ witnesses, out of reach for this representation.
Meaningful extension of kernel-checked range therefore runs through
structure rather than length: certifying Hughes's modular reduction
under the same axiom discipline would replace per-$n$ witnesses by
per-residue-class arguments, and the frontier search cells themselves
have exactly the shape of our witnesses. The two artifacts are
complementary today --- his reach with a mixed trust base, our trust
base on a short range --- and composing them is the natural next step,
though replacing \texttt{native\_decide} throughout a development of
that size is a project in its own right.

On the question itself the present work is silent beyond $10^8$, and
honesty requires repeating what the heuristics say: any solution, if
one exists at all, lies beyond the current $9.17 \times 10^{18}$
frontier, in ranges where the prime-tuple constraints price its
existence very low. The value of a certificate here is not the hunt but
the floor it establishes --- and the demonstration, on a problem where
one incorrect proof has already circulated, that the finite part of a
domination argument can be settled by a kernel rather than asserted.

\section{Precise claims}\label{sec:claims}

Certified, in Lean 4.30.0 / mathlib v4.30.0, axiom closure exactly
$\{\mathtt{propext}, \mathtt{Classical.choice}, \mathtt{Quot.sound}\}$,
enforced in CI \cite{decanus}:
no $n$ with $24 < n \le 10^7$, and no $n$ with $24 < n \le 10^8$,
satisfies the problem's inequality, stated against the pinned
\texttt{formal-conjectures} predicate.

Cited, computational, outside any kernel: no solution in
$(24, 10^{12}]$ (Id\'en \cite{iden}); no solution in $(24,\,
6.16 \times 10^{17}]$, extended to $(24,\, 9.17 \times 10^{18}]$,
within the Hughes reduction whose Lean layer uses
\texttt{native\_decide} plus one problem-specific axiom
(Hughes \cite{hughesrepo}, bentrd \cite{bentrd}).

We claim nothing else. The structural mathematics of the problem
belongs to the discussion-thread authors cited above; the search
frontier belongs to Id\'en, Hughes, and bentrd.

\subsection*{Acknowledgements}
Claude (Anthropic) was used throughout the development of the Lean
proofs, the generation pipeline, and this note. Every result is
kernel-checked, and the pipeline is cross-verified by the independent
implementations of Section~\ref{sec:validation}.

\begin{thebibliography}{99}

\bibitem{erdos79}
P.~Erd\H{o}s,
\emph{Some unconventional problems in number theory},
Math.\ Mag.\ \textbf{52} (1979), 67--70.

\bibitem{guy}
R.~K.~Guy,
\emph{Unsolved Problems in Number Theory},
3rd ed., Springer, 2004.

\bibitem{bloom647}
T.~F.~Bloom,
\emph{Erd\H{o}s Problem \#647},
\url{https://www.erdosproblems.com/647}, and the associated discussion
thread, \url{https://www.erdosproblems.com/forum/discuss/647}.
Accessed 17 August 2026.

\bibitem{oeisA087280}
OEIS Foundation Inc.,
\emph{Sequence A087280},
\url{https://oeis.org/A087280}.

\bibitem{oeisA062249}
OEIS Foundation Inc.,
\emph{Sequence A062249},
\url{https://oeis.org/A062249}.

\bibitem{iden}
P.~Id\'en,
\emph{Computational verification of Erd\H{o}s problem 647 up to
$10^{12}$: gap-growth analysis, depth-record extension to $D=11$, and
structural analysis}, v3, Zenodo, 2026.
\url{https://doi.org/10.5281/zenodo.21084248}.

\bibitem{hughesrepo}
S.~D.~Hughes,
\emph{erdos647-proof-chain: Lean 4 proof-chain package for Erd\H{o}s
problem \#647}, with the frontier certificate,
\url{https://github.com/scottdhughes/erdos647-proof-chain}, 2026.

\bibitem{bentrd}
bentrd,
\emph{Erd\H{o}s \#647 frontier extension to $9.174 \times 10^{18}$},
\url{https://github.com/bentrd/erdos647-frontier-extension}, 2026.

\bibitem{agbanwa}
J.~Agbanwa,
\emph{AI assisted (possible) solution to Erd\H{o}s problem 647},
Zenodo, 2026. \url{https://zenodo.org/records/18390414}.
See the assessment in the discussion thread of \cite{bloom647}.

\bibitem{formalconjectures}
Google DeepMind,
\emph{formal-conjectures},
\texttt{FormalConjectures/ErdosProblems/647.lean} at commit
\texttt{c252a410},
\url{https://github.com/google-deepmind/formal-conjectures}.

\bibitem{lean4}
L.~de~Moura and S.~Ullrich,
\emph{The Lean 4 theorem prover and programming language},
CADE 28, Springer, 2021, 625--635.

\bibitem{mathlib}
The mathlib Community,
\emph{The Lean mathematical library},
CPP 2020, ACM, 2020, 367--381.

\bibitem{lean4checker}
The Lean prover community,
\emph{lean4checker},
\url{https://github.com/leanprover/lean4checker}.

\bibitem{decanus}
I.~Mian and S.~Siddique,
\emph{Decanus: a kernel-checked exclusion certificate for Erd\H{o}s
problem 647}, v1.0.0,
\url{https://github.com/ibrahimmian36/decanus}, 2026.

\end{thebibliography}

\end{document}
